\documentclass{article}

% E&D Track, double-blind (default for the track)
\usepackage[eandd]{neurips_2026}

\usepackage[utf8]{inputenc}
\usepackage[T1]{fontenc}
\usepackage{hyperref}
\usepackage{url}
\usepackage{booktabs}
\usepackage{amsfonts}
\usepackage{amsmath}
\usepackage{amssymb}
\usepackage{nicefrac}
\usepackage{microtype}
\usepackage{xcolor}
\usepackage{graphicx}

\DeclareMathOperator*{\argmax}{arg\,max}

\newcommand{\dataset}{\textsc{CTSpinoPelvic1K}}
\newcommand{\TS}{TotalSegmentator}
\newcommand{\SOTA}{TotalSegmentator}
\newcommand{\sota}{TS}
\newcommand{\VERIDAH}{VERIDAH}
\newcommand{\VerSe}{VerSe}

\title{Not All Spines Are Created Equal: How CT Segmenters\\
Fail on Lumbosacral Transitional Vertebrae}

\author{%
  Anonymous Author(s) \\
  Anonymous Affiliation \\
  \texttt{anonymous@institution.edu}
}

\begin{document}

\maketitle

% ============================================================================
\begin{abstract}
Wrong-level spine surgery remains a persistent
never-event, with the Lumbosacral Transitional Vertebra
(LSTV, 5--35\% prevalence) as its largest anatomic
driver. We audit \TS{}---the most widely deployed CT
segmentation system and default back-end of
3D~Slicer---and show that on lumbarization cases its
output is anatomically indistinguishable from a
wrong-level surgical plan: lacking an L6 class, it
shifts lumbar labels caudally by one level, with
junction-DSC dropping 29 points on Any-LSTV. No
existing CT benchmark surfaces this---not \TS{}'s own
evaluation, not \citet{Li2025LumbarMechanism} (the
per-class DSC frontier on the matched COLONOG split),
and not \VERIDAH{} (vertebra-\emph{labeling} SOTA on a
private CT cohort, operating on pre-localized crops
with no pelvis). We
contribute: (i)~an LSTV-stratified evaluation protocol
(per-class DSC, junction-DSC over a 40\,mm L5/S1
window, per-class voxel confusion at the junction)
translating the level-shift mechanism into quantifiable
wrong-level surgical risk;
(ii)~\dataset{}---1{,}153 CT volumes across 802
patients with unified spinopelvic masks, 33
LSTV-positive cases across a 6-way phenotype taxonomy,
and radiologist Castellvi typing (I--IV) on
all 33; (iii)~the first publicly-released LSTV-aware
dual-spinopelvic CT segmenter, exceeding \TS{} on
sacrum (\TS{}=0.817 zero-shot) and uniquely delivering
both spine and pelvis in a single forward pass with
released weights; (iv)~an empirical isolation showing
\VERIDAH{}'s training-side L5/L6 merge eliminates the
collision driving \TS{}'s level shift but exposes a
residual L4/last\_lumbar collision on count-style
sacralization, motivating \VERIDAH{}'s sequence
predictor as the necessary disambiguation component.
\dataset{}, the protocol, and 5-fold checkpoints are
released publicly.
\end{abstract}

% ============================================================================
\section{Introduction}
\label{sec:intro}

Wrong-level spine surgery is a persistent never-event in
operative spine
care~\citep{Konin2010LumbosacralRelevance}. Despite
intraoperative imaging and surgical-pause checklists, the
most common cause---failure to correctly count vertebral
levels in the presence of an LSTV---continues to account for
the majority of reported wrong-level lumbar
incidents~\citep{Carrino2011EffectImaging}. The mechanism is
well documented: a surgeon counts lumbar bodies upward from
the sacrum, finds five (because lumbarization presents a
sixth lumbar body that has not been recognized), and labels
the bottom body as L5 when it is anatomically L6. Surgery
proceeds one level caudal to the intended target.

Modern radiology workflows increasingly route automated
segmentation as an upstream substrate for clinical
decision-making. The most widely deployed CT segmentation
system in this role is
\TS{}~\citep{Wasserthal2023TotalSegmentator:Images}: an
nnU-Net trained on 1{,}204 routine CTs to delineate 104
anatomic structures, the default segmentation back-end of
3D~Slicer, with thousands of clinical and research
installations. The original paper reports an aggregate
test-set Dice of 0.943 across all 104 classes and
qualitatively documents two spine failure modes
affecting neighboring-vertebrae label confusion (1\%
of cases) and neighboring-rib confusion (12\%). \TS{} is neither the
only public CT spine segmenter nor the per-class DSC
frontier on the matched CTSpine1K~+~CTPelvic1K (COLONOG)
split (\S\ref{sec:related}); it is, however, the model
surgeons and downstream tools consume. This paper
concerns one anatomic class on which \TS{} does not
perform at its headline level, and the consequences
when it does not.

\paragraph{The clinical claim.}
On lumbarized spines, \SOTA{} produces an output pattern
that is, when consumed by surgical planning pipelines or by
surgeons confirming operative levels, \emph{indistinguishable
from a wrong-level surgical plan}. Its label vocabulary
contains L1--L5 but no L6 class. On a lumbarized spine (six
lumbar bodies plus sacrum), the model distributes its five
lumbar labels across the six bodies by shifting caudally by
one anatomic level: the body that is anatomically L6
receives the L5 label, the body that is anatomically L5
receives the L4 label, and so on. A surgeon targeting the
disc between predicted-L4 and predicted-L5 lands on the
anatomical L5--L6 disc, one level caudal to the intended
L4--L5 target. \SOTA{} reproduces the canonical
LSTV-induced miscount automatically and silently.

\paragraph{The architectural alternative: \VERIDAH{}.}
Recent work~\citep{Moller2026VERIDAH:Sequences} addresses
transitional vertebra labeling via a two-stage
architecture: a 3D DenseNet169 classifier on per-vertebra
crops (four heads: label, region, transition, visibility),
followed by a constrained sequence predictor that imposes
anatomical-count constraints (4--6 lumbar bodies; the
predicted sequence may end in two consecutive L5 labels,
the second interpreted as L6). The classifier is trained
with L5/L6 (and T12/T13) labels merged for class-imbalance
reasons; disambiguation is deferred to the sequence
predictor. The architectural commitment is empirically
validated (97.22\% CT subject-level agreement) but has
never been tested against direct semantic segmentation,
and it is unknown whether a sufficiently expressive
voxel-wise network trained at full resolution with the
same merge requires the sequence predictor at all. We
supply the missing test.

\paragraph{The evaluation gap.}
Neither \SOTA{}, \VerSe{}~\citep{Sekuboyina2021VerSe:Images},
CTSpine1K~\citep{Deng2021CTSpine1K:Tomography}, nor
CTPelvic1K~\citep{Liu2021DeepModels} report LSTV-stratified
performance. Earlier CT labelers explicitly excluded the
failure mode~\citep{Sekuboyina2020LabelingAnatomy}.
\VERIDAH{} reports anomaly performance on a private cohort
without isolating semantic-segmentation contribution. The
result is a measurement vacuum: a production model emits a
confident, locally well-formed, caudally-shifted prediction
on a 5--35\%-prevalence congenital variant with no error
signal at evaluation time.

\paragraph{Closing the gap requires hardware.}
Fitting both lumbar spine and pelvis at
$0.75{\times}0.75{\times}0.80$\,mm spacing and
$256{\times}320{\times}320$ patch consumes ${\sim}100$\,GB
GPU memory under \texttt{nnUNetResEncUNetPlans\_100G},
available on H200 (140\,GB HBM3e) but not on the
L40S/A100/H100 GPUs that bracket \SOTA{}'s training
configuration. Neither \VERIDAH{} nor SPINEPS train at this
combined spinopelvic resolution.

\paragraph{Contributions.}
\begin{enumerate}
\item \textbf{Wrong-level surgical risk in \sota{}.}
  Systematic one-level caudal shift on lumbarization
  quantified through per-class DSC collapse (L4=0.388,
  L5=0.291), per-class voxel confusion at L5/S1
  (95\%+ of GT-L6 voxels labeled L5 on representative
  cases, \S\ref{sec:case_study}), and structural absence of
  L6 prediction across all 1{,}152 evaluation cases.
\item \textbf{\dataset{}: a stratified spinopelvic
  benchmark.} 1{,}153 abdominal CT records across 802
  patients, unified spinopelvic masks, 33 LSTV-positive
  patients across a 6-way subtype taxonomy
  (lumb $n=14$; sacr\_count $n=9$; full sacralization
  $n=6$; semi-sacralization $n=2$; ambiguous $n=2$), and a
  registration-free fusion pipeline that resolves a
  previously undocumented Y-axis coordinate-convention
  mismatch between CTSpine1K and CTPelvic1K. Both source
  datasets share the COLONOG cohort that underlies
  \VerSe{} and CTSpine1K, but neither encodes pelvic
  context with LSTV stratification.
\item \textbf{First LSTV-aware spinopelvic semantic
  segmenter.} A 5-fold CV nnU-Net 3d\_fullres baseline on
  H200 hardware adopting \VERIDAH{}'s training-side merge
  (last\_lumbar = L5+L6) at the voxel level, combined with
  pelvic-class supervision in a single forward pass.
  \VERIDAH{} and SPINEPS operate on the spine alone, and
  CTPelvic1K-trained pelvic segmenters are not LSTV-aware.
\item \textbf{State-of-the-art for openly-released CT
  dual spinopelvic segmentation.} Our segmenter is the
  only publicly-released model that delivers both spine
  vertebrae (with explicit LSTV handling via
  \texttt{last\_lumbar}) and pelvis (sacrum + hip) in a
  single forward pass; we exceed \TS{} on sacrum and
  establish best-in-class for the dual task
  (Table~\ref{tab:comparison}). \S\ref{sec:related}
  details the reproducibility tiers separating openly-released
  from paper-only baselines.
\item \textbf{Empirical isolation: merge alone is
  insufficient.} The merge eliminates the L5/L6 collision
  driving \sota{}'s level shift but exposes a residual
  L4/last\_lumbar collision on count-style sacralization.
  This residual collision empirically isolates the
  component of \VERIDAH{}'s success attributable to its
  sequence predictor---a direction \VERIDAH{}'s authors
  could not ablate.
\end{enumerate}

% ============================================================================
\section{Related Work}
\label{sec:related}

\paragraph{Wrong-level spine surgery and LSTV.}
\citep{Konin2010LumbosacralRelevance,Carrino2011EffectImaging,Sekharappa2014LumbosacralSignificance}
document LSTV-induced miscounting (5--35\% prevalence) but
do not address automated segmentation as a propagation
pathway for the miscount.
\citet{Kwak2025AutomatedModel} detect LSTV on plain
radiographs (binary task, different modality).

\paragraph{Production CT spine segmentation: \TS{}.}
\TS{}~\citep{Wasserthal2023TotalSegmentator:Images} is the
most widely deployed CT segmentation model in clinical and
research workflows---an nnU-Net trained on 1{,}204 routine
CTs, the default segmentation back-end of 3D~Slicer. Its
label vocabulary covers C1--L5 plus thoracic and cervical
levels but contains no L6 class and no
transitional-vertebra phenotype, and its training cohort
is not LSTV-stratified. We audit \TS{} not because it is
the per-class DSC frontier---it is not---but because it
is the production system that downstream surgical-planning
pipelines actually consume.

\paragraph{Reproducibility-tiered prior work.}
\citet{Li2025LumbarMechanism} hold the published
per-class DSC frontier on the matched COLONOG split
(731 matched scans, train/val/test 440/146/145), with
a cross-scale feature fusion + linear-attention
transformer beating nnU-Net by 1--3 DSC points on
lumbar classes. \emph{Neither weights nor inference
code are released}, making the result paper-only; like
\TS{}, they do not stratify by LSTV.
Table~\ref{tab:comparison} reproduces all seven of
their methods to anchor our pelvic claims.
\citet{Liu2025AutomaticDataset} release weights for
FracSegNet, a cascaded 3D nnU-Net pelvic segmenter
(PENGWIN + CTPelvic1K) providing sacrum and L/R hipbone
but no spine vertebrae; the authors note qualitative
degradation on sacralization cases when training on
CTPelvic1K alone (their \S4.3) but do not
quantify---corroborating evidence for the LSTV gap our
benchmark surfaces.

\paragraph{Vertebra-labeling SOTA: \VERIDAH{}.}
\citet{Moller2026VERIDAH:Sequences} address
transitional-vertebra \emph{labeling}, a task distinct
from segmentation. Given vertebra centroids supplied by an
upstream segmenter, a 3D DenseNet169 classifier with four
heads (label, region, transition, visibility) operates on
per-vertebra crops and a constrained sequence predictor
imposes anatomical-count constraints (4--6 lumbar
bodies; the predicted sequence may end in two
consecutive L5 labels, the second interpreted as L6).
On a private
in-house cohort of 1{,}536 CT volumes, \VERIDAH{} reports
99.18\% subject-level perfect-labeling agreement,
substantially above the prior CT-labeling
SOTA~\citep{Meng2023VertebraeCycle} (77.81\%). Three
properties of \VERIDAH{}'s setup matter for our framing.
\textbf{(1)~It does not segment}: it requires upstream
centroid localization, which the authors source from
\VerSe{}-derived~\citep{Sekuboyina2021VerSe:Images} models
or SPINEPS~\citep{Moller2024SPINEPSautomaticSegmentation};
SPINEPS itself is T2w MRI-only.
\textbf{(2)~It does not include the pelvis}; label space
is C1--L6 only.
\textbf{(3)~CT validation is on a private cohort},
so \VERIDAH{}'s CT performance is not reproducible on
public benchmarks. Its ablation (their Table~5) shows the
four-head + sequence-predictor configuration substantially
outperforms either component alone, but the comparison is
never made against direct semantic segmentation on a
public LSTV-stratified split. Our \dataset{}-trained
segmenter (\S\ref{sec:baseline}) supplies that missing
comparison.

\paragraph{Benchmarks and evaluation rigor.}
The \VerSe{}~\citep{Sekuboyina2021VerSe:Images} challenge
(460 CT volumes) is the only public CT benchmark whose
label vocabulary explicitly accommodates L6 (label 25),
but it is unstratified by LSTV subtype and contains no
pelvic ground truth. Earlier CT labelers explicitly
excluded the failure mode~\citep{Sekuboyina2020LabelingAnatomy}.
The position that benchmarks should report stratified
safety-critical
performance~\citep{Maier-Hein2024MetricsValidation}
motivates the protocol of \S\ref{sec:protocol}.

% ============================================================================
\section{Failure-Mode Analysis: \sota{} on LSTV}
\label{sec:audit}

We evaluate \SOTA{} v2 (\texttt{total} task) zero-shot at
full resolution on the full \dataset{} matched cohort.
\SOTA{} was not trained on any \dataset{} case, so every
case is valid; this provides $N{=}1{,}152$ successful
records aggregated to 801 patients via patient-level
merging, including 33 LSTV-positive patients.
Full-resolution mode isolates the failure as a property of
label vocabulary and training distribution rather than
inference resolution.

Table~\ref{tab:sota_stratified} reports performance by
phenotype. On normal anatomy ($n=768$) \SOTA{} performs at
its published level. On Any-LSTV ($n=33$) every metric
drops; lumbarization collapses caudally
(L1=0.878\,$\to$\,L4=0.388\,$\to$\,L5=0.291) with
junction-DSC at 0.357. Count-style sacralization ($n=9$)
collapses uniformly across all lumbar classes (L1=0.265
down to L5=0.000). \textsc{L6} predictions are
structurally zero across all 1{,}152 cases.

\begin{table}[t]
\caption{\sota{} zero-shot stratified by LSTV phenotype on
\dataset{}, full-resolution mode. Junction DSC over a
40\,mm L5/S1 axial window. \textsc{L6} count: cases with
any L6 voxel emitted (structurally zero by label
vocabulary). $\dagger$: junction analysis not computed for
separate-mode records.}
\label{tab:sota_stratified}
\centering
\small
\begin{tabular}{lrcccccccr}
\toprule
\bfseries Subgroup & \bfseries $n$ &
\bfseries L1 & \bfseries L2 & \bfseries L3 & \bfseries L4 & \bfseries L5 &
\bfseries Sac & \bfseries Hips &
\bfseries Jxn DSC \\
\midrule
All cases                              & 801 & 0.911 & 0.906 & 0.913 & 0.914 & 0.913 & 0.817 & 0.929 & 0.818 \\
Normal                                 & 768 & 0.921 & 0.918 & 0.929 & 0.936 & 0.938 & 0.817 & 0.929 & 0.824 \\
\midrule
Any LSTV                               & 33  & 0.681 & 0.615 & 0.523 & 0.367 & 0.314 & 0.822 & 0.929 & 0.534 \\
\quad Lumbarization                    & 14  & 0.878 & 0.853 & 0.690 & 0.388 & 0.291 & 0.833 & 0.928 & \textbf{0.357} \\
\quad Sacralization (count-style)      &  9  & 0.265 & 0.100 & 0.055 & 0.001 & 0.000 & 0.814 & 0.926 & 0.697 \\
\quad Semi-sacralization               &  2  & 0.951 & 0.951 & 0.951 & 0.955 & 0.956 & 0.824 & 0.937 & $\dagger$ \\
\quad Sacralization (full)             &  6  & 0.718 & 0.676 & 0.673 & 0.724 & 0.793 & 0.826 & 0.932 & 0.636 \\
\quad Ambiguous                        &  2  & 0.952 & 0.951 & 0.882 & 0.678 & 0.374 & 0.771 & 0.932 & $\dagger$ \\
\midrule
\textsc{L6 emit}                       &     &       &       &       &       &       &       &       & 0 / 1152 \\
\bottomrule
\end{tabular}
\end{table}

\paragraph{The level-shift mechanism.}
The asymmetry between lumbarization (severe, caudally
progressive) and count-style sacralization (uniform
collapse) follows from a single fact: \SOTA{}'s output head
has no L6 class. For 6-vertebra anatomies the model
compresses 6 anatomic units into 5 labels by shifting label
assignment caudally; the most rostral body remains L1
(escape into thoracic vertebrae), but every successive body
is misaligned, with the final body labeled L5 when it is
anatomically L6. For 4-vertebra anatomies the model emits
standard L1--L5, overshooting the GT lumbar count and
producing a $+1$-level offset uniformly. The level-shift
mechanism is the segmentation analog of the surgical
miscount~\citep{Konin2010LumbosacralRelevance}.

% ============================================================================
\section{Case Study: Wrong-Level Surgical Risk in \sota{}
Output}
\label{sec:case_study}

Aggregate metrics establish the population pattern; the
case study below establishes the mechanism at the level a
surgeon would consume the output. Token~167 is a fused-mode
lumbarization case with junction-DSC 0.221, the lowest
among lumbarization cases. Table~\ref{tab:case167} shows
the 40\,mm-window confusion matrix: 92{,}537 of 92{,}629
GT-L5 voxels (99.9\%) are predicted as L4, and 81{,}088 of
85{,}186 GT-L6 voxels (95.2\%) are predicted as L5---the
level-shift mechanism in its cleanest form.

\begin{table}[h]
\caption{Token~167 (lumbarization) per-class voxel
confusion matrix over the 40\,mm L5/S1 window. Bold cells:
dominant prediction class for each GT row. The
diagonal-shifted dominance pattern (GT-L5 $\to$ pred L4;
GT-L6 $\to$ pred L5) is the empirical signature of the
level-shift failure.}
\label{tab:case167}
\centering
\small
\begin{tabular}{lrrrrrr}
\toprule
\bfseries GT $\backslash$ Pred & \bfseries bg & \bfseries L3 & \bfseries L4 & \bfseries L5 & \bfseries sacrum & \bfseries L6 \\
\midrule
L5     & 4{,}572  & 2 & \textbf{92{,}537} & 90              & 0           & 0 (n/a) \\
L6     & 4{,}089  & 0 & 9                 & \textbf{81{,}088} & 0         & 0 (n/a) \\
sacrum & 68{,}635 & 0 & 0                 & 36              & \textbf{80{,}334} & 0 (n/a) \\
\bottomrule
\end{tabular}
\end{table}

\paragraph{Surgical translation.}
A surgeon planning an L4--L5 microdiscectomy and using the
segmentation to confirm operative level identifies the
disc between predicted-L4 (anatomical L5) and predicted-L5
(anatomical L6). The targeted disc is L5--L6, one level
caudal to the L4--L5 intent. The output is locally
well-formed: five compact vertebral predictions in roughly
correct positions. The error is detectable only by
counting bodies (six lumbar, not five), which downstream
pipelines do not perform.

% ============================================================================
\section{\dataset{}: Construction and Composition}
\label{sec:dataset}

\dataset{} fuses CT scans from
COLONOG~\citep{Smith2015DataCOLONOGRAPHY} with annotations
from CTSpine1K~\citep{Deng2021CTSpine1K:Tomography}
(\VerSe{}-derived per-vertebra masks) and
CTPelvic1K~\citep{Liu2021DeepModels} (sacrum and bilateral
hips). The same COLONOG cohort underlies \VerSe{} and
CTSpine1K; \dataset{}'s contribution is the patient-level
fusion of CTSpine1K's vertebra masks with CTPelvic1K's
pelvic masks plus the LSTV-stratification overlay neither
source provides. Of 1{,}153 matched records, 342 carry
fused (spine + pelvis) ground truth, 352 are separate-mode,
90 are spine-only, and 20 are pelvic-only
(Table~\ref{tab:scan_stats}).

\begin{table}[t]
\caption{\dataset{} composition; LSTV phenotype
breakdown in Table~\ref{tab:castellvi_dist}.}
\label{tab:scan_stats}
\centering
\small
\begin{tabular}{lcc}
\toprule
\bfseries Property & \bfseries Value & \bfseries LSTV cases \\
\midrule
Patients matched (\dataset{})           & 802     & --- \\
Records matched (\dataset{})            & 1{,}153 & 33 patients (44 records) \\
\quad Fused (spine + pelvis)            & 342     & --- \\
\quad Separate (spine + pelvic)         & 352     & --- \\
\quad Spine-only / Pelvic-only          & 90 / 20 & --- \\
\bottomrule
\end{tabular}
\end{table}

\paragraph{Patient-anchored fusion and coordinate
correction.} Patient identity is resolved via canonical UID
matching (825/825 patients). Each mask is assigned to its
source CT acquisition by intra-patient bone-HU
maximization,
$s^*(m) = \argmax_{s \in \mathcal{S}_p}
|\{v \in \mathcal{L}(m, s) : I_s(v) > 200\,\text{HU}\}|/|\mathcal{L}(m, s)|$,
a metadata-free oracle since \texttt{SeriesDescription}
fields are unreliable in TCIA releases (99.6\% spine,
100\% pelvic placement). CTSpine1K and CTPelvic1K were
converted from DICOM with different \texttt{dcm2niix}
invocations, producing a systematic Y-axis sign reflection
in NIfTI affines; we detect via affine-derived center sign
and reverse along axis~1 with affine update. We adopt
\VERIDAH{}'s training-side merged-label scheme at the
voxel level: 0\,=\,bg, 1--4\,=\,L1--L4,
\textbf{5\,=\,last\_lumbar} (anatomical L5 in 5-vertebra
spines; both L5 and L6 in 6-vertebra spines), 6\,=\,sacrum,
7--8\,=\,left/right hip, 9\,=\,ignore.

\paragraph{Six-way LSTV taxonomy with dual-lineage
labels.} We retain both LSTV label sources separately and
partition into six subtypes by combining filename-qualifier
evidence (CTPelvic1K morphological criterion) with
vertebral-counting evidence (CTSpine1K
\VerSe{}-label-presence heuristic):
\textbf{normal}, \textbf{lumb} (6 vertebrae labeled or
\VerSe{} label 25 present), \textbf{sacr\_count} (4
vertebrae labeled with sacralization flag),
\textbf{semisacralization}, \textbf{sacralization}
(both via filename qualifier), and \textbf{ambiguous}
(cross-source disagreement). An earlier substring-match
parser bug collapsed semisacralization into full
sacralization---merging two anatomically distinct
phenotypes whose failure modes diverge sharply (full
sacralization preserves the lumbar count and produces
near-normal per-class DSC, whereas semisacralization
introduces a partial transition that confuses the L5/S1
boundary). Any benchmark reporting an aggregate
sacralization DSC under this collapse averages over
two qualitatively different model behaviors; we identify
and fix this. Without the corrected 6-way taxonomy, the
qualitatively different failure patterns across LSTV
variants (\S\ref{sec:audit}) are invisible.

\paragraph{Partial-annotation training protocol.}
\label{sec:partial}
Training only on fused records discards 70\% of the data.
We use nnU-Net v2's
\texttt{ignore\_label}~\citep{Isensee2024Nnu-netSegmentation}
to train on all 1{,}153: voxels in unsupervised regions are
written as label~9, contributing zero to Dice and
cross-entropy. Spine classes (1--5) receive supervision
from 784 records; pelvic classes (6--8) from 714. This is a
$3.36\times$ effective expansion. Per-record subtype
refinement gates training-time tags on supervised anatomy.

\paragraph{Castellvi annotation layer.}
\label{sec:castellvi}
The 6-way taxonomy captures the \emph{count} dimension of
LSTV. The Castellvi
classification~\citep{Castellvi1984LumbosacralDefects}
captures the orthogonal \emph{morphological} axis:
Type~I (dysplastic transverse process), II
(pseudoarticulation), III (bony fusion), or IV (mixed
II/III), each \texttt{a}~(unilateral) or
\texttt{b}~(bilateral). It is the surgeon's
pre-operative wrong-level-risk indicator, and what a CT
segmenter would need to predict to flag LSTV for review
rather than silently emit a caudally-shifted label.
Radiologists annotated all 33 LSTV cases
(extending the schema, originally sacralization-specific,
to score the bony junction uniformly across variants),
with per-side \texttt{a}/\texttt{b} laterality.

\begin{table}[h]
\caption{Castellvi distribution across all 33 LSTV cases
of \dataset{}, cross-tabulated by 6-way count taxonomy.
Reads by radiologists.}
\label{tab:castellvi_dist}
\centering
\small
\setlength{\tabcolsep}{4pt}
\begin{tabular}{lrrrrrrrr}
\toprule
\bfseries Subtype & \bfseries Ia & \bfseries Ib & \bfseries IIa & \bfseries IIb & \bfseries IIIa & \bfseries IIIb & \bfseries IV & \bfseries Total \\
\midrule
Lumbarization                  & 0 & 1 & 1 & 3 & 1 & 7 & 1 & 14 \\
Sacralization (count-style)    & 0 & 0 & 0 & 0 & 0 & 9 & 0 & 9 \\
Sacralization (full)           & 0 & 0 & 0 & 0 & 1 & 3 & 2 & 6 \\
Semi-sacralization             & 0 & 0 & 1 & 0 & 1 & 0 & 0 & 2 \\
Ambiguous                      & 0 & 0 & 0 & 1 & 0 & 0 & 1 & 2 \\
\midrule
Total                          & 0 & 1 & 2 & 4 & 3 & 19 & 4 & 33 \\
\bottomrule
\end{tabular}
\end{table}

\paragraph{Morphology-stratified failure interpretation.}
The cross-tabulation (Table~\ref{tab:castellvi_dist})
surfaces two findings with direct bearing on model
behavior. (1)~\textbf{All 9 count-style sacralization
cases are Castellvi~IIIb}---bilateral bony fusion---and
19/33 (58\%) of LSTV overall is IIIb. The residual
L4/last\_lumbar collision (\S\ref{sec:residual})
therefore occurs precisely when the discriminating cue
is locally observable: bony fusion of the bottom lumbar
to sacrum is visually distinct from a normal L5--S1
junction, yet the segmenter still emits last\_lumbar
for those voxels. Direct semantic segmentation has not
learned the Castellvi correlate of the count anomaly
even when it is morphologically present in the input---an
empirical argument for downstream Castellvi-aware
reasoning at the lumbosacral junction (whether
VERIDAH-style global count constraints or a dedicated
Type~III morphology head). (2)~Only 1/33 cases is Type~I
(population prevalence
14--16\%~\citep{Sekharappa2014LumbosacralSignificance}),
reflecting source-dataset selection bias toward
morphologically conspicuous presentations. Reported
performance therefore characterizes the conspicuous
regime; the clinically subtler Type~I fraction is
uncharacterized in this benchmark.

% ============================================================================
\section{Evaluation Protocol}
\label{sec:protocol}

We define four metrics as the minimum required to attribute
LSTV failure modes to wrong-level surgical
risk~\citep{Maier-Hein2024MetricsValidation}: (1) per-class
DSC stratified by 6-way LSTV subgroup, exposing class-level
degradation; (2) junction-DSC over a 40\,mm L5/S1 axial
window, the clinical window for surgical planning;
(3) per-class voxel confusion at the L5/S1 junction, which
exposes structured failure modes (the level shift) that
DSC alone cannot characterize; and (4) for the merged-label
baseline, three focused per-epoch confusion blocks anchored
on the L4/last\_lumbar boundary that drives the residual
collision (\texttt{last\_lumbar\_GT\_on\_lumb},
\texttt{L4\_GT\_on\_sacr\_count},
\texttt{pred\_last\_lumbar\_on\_sacr\_count}) plus a
specificity number
(\texttt{last\_lumbar\_specificity\_on\_sacr\_count}).

% ============================================================================
\section{Spinopelvic Baseline}
\label{sec:baseline}

\paragraph{Architecture and training.}
nnU-Net v2~\citep{Isensee2020NnU-Net:Segmentation,Isensee2024Nnu-netSegmentation}
ResEnc-M planner, \texttt{3d\_fullres},
\texttt{nnUNetResEncUNetPlans\_100G} (100\,GB GPU budget).
Spacing $0.80{\times}0.75{\times}0.75$\,mm, patch
$\mathbf{256{\times}320{\times}320}$, batch~2; 7-stage
residual encoder
($n_\text{features}{=}[32,64,128,256,320,320,320]$,
$n_\text{blocks}{=}[1,3,4,6,6,6,6]$, $3^3$ kernels, $2^3$
strides, instance norm, deep supervision at 5 levels);
9-channel output (bg + 8 fg); ignore-label masked from
loss. Class-reweighted DC\_and\_CE loss
($\mathbf{w}{=}[0.5,1,1,1,1,\mathbf{1.5},\mathbf{2.0},1,1]$;
last\_lumbar boost targets the clinically-important
boundary, sacrum boost supports Castellvi precision); SGD
Nesterov 0.99, $\lambda{=}3{\times}10^{-5}$, PolyLR
$10^{-2}$, 500 epochs (250 train + 50 val iters).

\paragraph{Two-tier LSTV emphasis.}
\textbf{Queue-level oversampling}: all five non-normal
subtypes duplicated to 25\% of per-fold training samples
(matching \VERIDAH{}'s 4$\times$ rate).
\textbf{Class-reweighted CE} (above). No per-patch
foreground biasing: under merged labels, biasing toward
an L6 class is ill-defined, and biasing the
L4/last\_lumbar boundary on sacr\_count would confound
the residual-collision measurement.

\paragraph{Hardware.}
H200 NVL (140\,GB HBM3e), 24 CPU workers, preunpacking
once per job to a 1.7\,TB local-NVMe cache (replacing NFS
reads that throttled at 21\,MB/s under 5-fold contention).
Per-epoch ${\approx}\,400$--$700$\,s; per-fold
${\approx}\,3$ days for 500 epochs. Reducing the GPU memory
budget forces either smaller patch (loses spinopelvic
context) or coarser spacing (loses Castellvi-defining
detail); 100\,GB is the smallest budget fitting both.
5-fold patient-stratified CV with 6-way LSTV
stratification.

% ============================================================================
\section{Results}
\label{sec:results}

\textit{Result tables populated from completed 5-fold CV.}

\paragraph{Headline.}
The merged-label spinopelvic baseline closes the LSTV
junction-DSC gap on lumbarization by
$\Delta{=}\,$\textbf{[+XX.X]}~points
(\sota{}\,=\,0.357\,$\to$\,ours\,=\,\textbf{[0.XXX]}, mean
across 5 folds) and emits non-degenerate last\_lumbar on
\textbf{[XX/14]} lumbarization cases with mean per-case
last\_lumbar DSC = \textbf{[0.XXX]}, eliminating the
structural absence that drives \sota{}'s level shift. A
residual class collision persists on count-style
sacralization (per-case L4 DSC = \textbf{[0.XXX]});
\S\ref{sec:residual} characterizes this as the empirical
signature of direct semantic segmentation's limit on
transitional anatomy. On pelvic classes, our model
exceeds \TS{} on sacrum by \textbf{[+XX.X]}~points and
establishes state-of-the-art for openly-released dual
spinopelvic CT segmentation (Table~\ref{tab:comparison}).

\subsection{Comparison to Prior CTSpine1K/CTPelvic1K
Baselines}
\label{sec:comparison}

Table~\ref{tab:comparison} reports per-class DSC across
three reproducibility tiers (paper-only, reproducible
baselines zero-shot on \dataset{}, and ours).

\begin{table}[t]
\caption{Per-class DSC ($\times 100$). Block-level
reproducibility tiers labeled in-table. The L5* column
is unmerged L5 for upper-block rows and merged
\textbf{last\_lumbar} (L5 in 5-vertebra spines, L5+L6
in 6-vertebra) for ours. Hip$_\text{m}$ is the
merged-hip class (LH+RH on both prediction and GT);
per-side LH and RH columns are dashed for our rows
since we train and evaluate on Hip$_\text{m}$ only.
\TS{}'s L5* is dashed (no merged class); FracSegNet
spine entries are dashed (pelvis-only). Bold: best in
column \emph{among reproducible models}.}
\label{tab:comparison}
\centering
\footnotesize
\setlength{\tabcolsep}{4pt}
\begin{tabular}{lccccccccc}
\toprule
\bfseries Method & \bfseries L1 & \bfseries L2 &
\bfseries L3 & \bfseries L4 & \bfseries L5* &
\bfseries Sacrum & \bfseries LH & \bfseries RH &
\bfseries Hip$_\text{m}$ \\
\midrule
\multicolumn{10}{l}{\emph{Reproduced from \citet{Li2025LumbarMechanism}, Tables~2--3 (paper-only; no released weights or code)}} \\
3D U-Net                & 77.19 & 78.34 & 81.71 & 82.86 & 92.13 & 90.02 & 93.74 & 93.42 & --- \\
ResNet-50               & 80.29 & 76.41 & 80.62 & 83.77 & 91.05 & 88.37 & 94.72 & 94.01 & --- \\
ResNet-101              & 74.57 & 78.11 & 78.54 & 83.84 & 86.35 & 89.05 & \textbf{95.69} & 94.42 & --- \\
SwinUNETR               & 88.87 & 83.12 & 84.36 & 89.04 & 92.15 & 91.38 & 92.33 & 93.38 & --- \\
AttentionUNet           & 80.26 & 75.84 & 79.64 & 82.32 & 90.39 & 92.48 & 93.41 & 95.45 & --- \\
nnU-Net~\citep{Isensee2020NnU-Net:Segmentation}     & 89.35 & 91.37 & 92.41 & 93.37 & 95.74 & 96.72 & 95.50 & \textbf{96.78} & --- \\
Li et al.\ (CFFM+LT)~\citep{Li2025LumbarMechanism}  & \textbf{92.98} & \textbf{92.36} & \textbf{93.92} & \textbf{93.31} & \textbf{97.07} & 96.98 & 94.37 & 96.77 & --- \\
\midrule
\multicolumn{10}{l}{\emph{Reproducible baselines (publicly-released weights), evaluated zero-shot on \dataset{}}} \\
\TS{}~\citep{Wasserthal2023TotalSegmentator:Images}  & [.XX] & [.XX] & [.XX] & [.XX] & ---   & [.XX] & [.XX] & [.XX] & [.XX] \\
FracSegNet~\citep{Liu2025AutomaticDataset}            & ---   & ---   & ---   & ---   & ---   & [.XX] & [.XX] & [.XX] & [.XX] \\
\midrule
\multicolumn{10}{l}{\emph{Ours: \dataset{} 5-fold CV, partial-annotation supervision, 6-way LSTV taxonomy}} \\
Ours, Normal ($n{=}$XX)   & [.XX] & [.XX] & [.XX] & [.XX] & [.XX] & \textbf{[.XX]} & --- & --- & \textbf{[.XX]} \\
Ours, Any-LSTV ($n{=}$33) & [.XX] & [.XX] & [.XX] & [.XX] & [.XX] & [.XX]          & --- & --- & [.XX] \\
\bottomrule
\end{tabular}
\end{table}

\paragraph{LSTV-stratified per-subgroup view.}
Table~\ref{tab:baseline} reports per-subgroup DSC and
junction-DSC for the same model, exposing the
LSTV-induced degradation that aggregate per-class DSC
masks. Compared to \sota{}'s zero-shot lumbarization
performance (Table~\ref{tab:sota_stratified}), our
last\_lumbar emission rate on lumb cases rises from
structural zero to \textbf{[XX/14]} and junction-DSC rises
from 0.357 to \textbf{[0.XXX]}, eliminating the level-shift
mechanism. The residual L4 DSC drop on sacr\_count cases
(\S\ref{sec:residual}) persists.

\begin{table}[t]
\caption{\dataset{} 5-fold CV baseline by LSTV subgroup,
held-out validation folds. \textsc{LL emit}: cases where
the model emitted $\geq 100$ last\_lumbar voxels.
$\Delta$ versus \sota{} zero-shot
(Table~\ref{tab:sota_stratified}).}
\label{tab:baseline}
\centering
\small
\begin{tabular}{lrcccccr}
\toprule
\bfseries Subgroup & \bfseries $n$ &
\bfseries L1--L4 & \bfseries last\_lumbar & \bfseries Sac &
\bfseries Hip$_\text{m}$ & \bfseries Jxn DSC & \bfseries LL emit \\
\midrule
Normal                  & XX & 0.XXX & 0.XXX & 0.XXX & 0.XXX & 0.XXX & XX/XX \\
Lumbarization           & 14 & 0.XXX & \textbf{0.XXX} & 0.XXX & 0.XXX & \textbf{0.XXX} & XX/14 \\
Sacralization (count)   &  9 & 0.XXX & ---   & 0.XXX & 0.XXX & 0.XXX & X/9   \\
Semi-sacralization      &  2 & 0.XXX & 0.XXX & 0.XXX & 0.XXX & 0.XXX & X/2   \\
Sacralization (full)    &  6 & 0.XXX & 0.XXX & 0.XXX & 0.XXX & 0.XXX & X/6   \\
Ambiguous               &  2 & 0.XXX & 0.XXX & 0.XXX & 0.XXX & 0.XXX & X/2   \\
\midrule
$\Delta$ vs.\ \sota{} (lumb)        & --- & +XX.X & N/A & +X.X & 0 & \textbf{+XX.X} & --- \\
$\Delta$ vs.\ \sota{} (sacr\_count) & --- & +XX.X & --- & +X.X & 0 & +X.X    & --- \\
\bottomrule
\end{tabular}
\end{table}

\paragraph{Hip class merging.}
\label{sec:hip_laterality}
We merge left and right hip into a single bone class
(Hip$_\text{m}$) on both prediction and ground truth
(Table~\ref{tab:comparison}) to avoid L/R class confusion
and provide a single hip-bone metric directly comparable
across models.

\subsection{The Residual Class Collision}
\label{sec:residual}

The merged scheme eliminates the L5/L6 collision driving
\sota{}'s wrong-level mechanism: GT-last\_lumbar voxels on
lumb cases are predicted as last\_lumbar in
\textbf{[XX]\%} of voxels (Block~1), against \sota{}'s
structurally-zero last\_lumbar emission. But it exposes a
residual collision at the L4/last\_lumbar boundary on
sacr\_count cases: GT-L4 voxels are predicted as
last\_lumbar in \textbf{[XX]\%} of voxels (Block~2), and
predicted-last\_lumbar voxels on sacr\_count cases are
anatomically L4 in \textbf{[XX]\%} (Block~3).
Specificity = \textbf{[0.XXX]}.

\begin{table}[t]
\caption{Residual class confusion under merged labels,
focused on L4/last\_lumbar. Block~1: GT-last\_lumbar on
lumb (good signal). Block~2: GT-L4 on sacr\_count
(failure-mode signal). Block~3: pred-last\_lumbar on
sacr\_count (inverse direction). [Numbers fill at training
completion.]}
\label{tab:confusion_blocks}
\centering
\small
\begin{tabular}{lcccc}
\toprule
\bfseries Block & last\_lumbar & sacrum & L4 & background \\
\midrule
Block~1: last\_lumbar GT on lumb        & \textbf{XX.X\%} & X.X\% & X.X\% & X.X\% \\
Block~2: L4 GT on sacr\_count           & \emph{X.X\%}    & X.X\% & \textbf{XX.X\%} & X.X\% \\
Block~3: pred-last\_lumbar on sacr\_count & ---           & X.X\% & \textbf{XX.X\%} & X.X\% \\
\bottomrule
\end{tabular}
\end{table}

\paragraph{Interpretation.}
The model has learned the spinopelvic spatial relationship
on normal anatomy and recovers the merged last\_lumbar
class on lumbarization. On count-style sacralization,
however, the positional prior dominates: the bottom-most
lumbar visible in a given patch is treated as
last\_lumbar regardless of the global vertebral count.
The disambiguating signal (vertebra count from a known
anchor) is non-local; a $256{\times}320{\times}320$
patch does not reliably contain both the caudal
lumbosacral and rostral cervico-thoracic boundaries
needed to count. This is a finding rather than a defect
of the baseline---it executes the cleanest test of
\VERIDAH{}'s training-side merge in the
semantic-segmentation setting and confirms the
architectural hypothesis underlying its two-stage
design: the merge addresses class imbalance, but
distinguishing the last-lumbar voxels of a 4-vertebra
spine from those of a 5-vertebra spine must be deferred
to a downstream component with access to global
anatomical structure.

% ============================================================================
\section{Discussion and Future Work}
\label{sec:discussion}

\paragraph{Generalization and hardware.}
The level-shift mislabeling failure is a property of any
spine segmentation model whose output head excludes L6,
deployed on a distribution containing lumbarization. The
100\,GB patch budget is presently only available on
H200-class hardware; smaller budgets will yield worse LSTV
gap closure and therefore residual wrong-level risk in any
downstream pipeline that consumes the output.

\paragraph{Future work.}
\textbf{(1)~\VerSe{} expansion via sacrum
pseudolabeling.} \VerSe{}~\citep{Sekuboyina2021VerSe:Images}
contains 460 CT volumes labeling L6 (label~25) but no
pelvis; our segmenter's sacral DSC (\textbf{[0.XX]})
supports sacrum pseudolabeling on \VerSe{} for a
$1{,}613$-volume spinopelvic cohort.
\textbf{(2)~Two-stage integration with \VERIDAH{}.}
Our segmenter can replace \VERIDAH{}'s \VerSe{}-derived
CT upstream, supplying sacrum and per-vertebra centroids
in a single forward pass.

\paragraph{Limitations.}
COLONOG is a US screening cohort (ages 45--75);
generalization to emergency/pediatric populations is not
evaluated. LSTV labels are inherited from CTPelvic1K and
CTSpine1K heuristics; Castellvi annotation covers all 33
LSTV cases but does not screen the \texttt{normal} cohort
for Type~I, and the source-dataset construction biases
against morphologically-mild Type~I presentations
(\S\ref{sec:castellvi}). Two count-axis subtypes
(semi-sacralization, ambiguous) have $n=2$. \sota{} is
evaluated zero-shot, the correct comparison for using
\sota{} as a production system but undersells what it
could achieve if retrained with an L6 class.

\paragraph{Implications.}
A production CT segmentation system that emits a confident,
locally well-formed, caudally-shifted prediction on a
5--35\%-prevalence congenital variant produces no error
signal at inference time and, when consumed downstream,
reproduces the canonical mechanism of wrong-level spine
surgery. Any future spine segmentation model claiming
readiness for surgical-planning integration should report
per-subgroup DSC and per-class confusion at the L5/S1
junction on a stratified LSTV cohort.

% ============================================================================
\section*{Data and Code Availability}

\dataset{} is hosted on HuggingFace Datasets (anonymous
repository for review) with both core and Responsible AI
fields of the Croissant
specification~\citep{Akhtar2024Croissant:Datasets}; a 50-record
sample covering all six LSTV subtypes is provided. The
evaluation protocol, data-construction pipeline, and
trainer source are released under permissive license at an
anonymous GitHub repository accessible without credentialed
access. Trained 5-fold checkpoints will be made available.

\bibliographystyle{plainnat}
\bibliography{bib/references}

\newpage
\input{checklist.tex}

\end{document}
